\section{Sharpness in the forced-midpoint profile model}
\label{pfobs:section}

This section concerns only the elementary scale-selection problem. Its examples
are not planar distance-set counterexamples. In particular, they do not show that
the sufficient Hausdorff--packing condition of this manuscript is necessary.

For a continuous function \(g:[0,1]\to\mathbb R\), an edge from \(n\) to
\(m<n\) has cost
$$
 c_g(m,n)=g(m)-\min_{[m,n]}g,
$$
and is admissible if \(2n-m\le1\). A forced-midpoint chain starts at
\(1-\zeta\), passes through \(1/2\), and ends at zero. Write
\(\mathcal C_\zeta(g)\) for the infimum of the total costs of these chains.
The example below attains this infimum with a finite chain.

\begin{proposition}[Exact two-collapse profile]\label{pfobs:exact}
Suppose
$$
 0<a\le\frac14,\qquad a\le b\le1,\qquad
 t:=\frac{1+b}{2(1+a)}\le\frac34,
$$
and
\begin{equation}\label{pfobs:barrier-condition}
 2(1-b)(1+a)^2\le(1+b)^2(1-a).
\end{equation}
Put
$$
 m=at,\qquad v=\frac b2,\qquad
 p=\frac{1+a+t-m}{2},\qquad
 C=\frac{1-t+m-a}{3}=\frac{(1-a)(1-t)}3.
$$
Define
\begin{equation}\label{pfobs:profile}
 g(x)=\begin{cases}
 bx,&0\le x\le1/2,\\
 v+1/2-x,&1/2\le x\le t,\\
 m+x-t,&t\le x\le p,\\
 a+1-x,&p\le x\le1.
 \end{cases}
\end{equation}
Then \(g\) is 1-Lipschitz and satisfies \(ax\le g(x)\le bx\).
For every \(0<\zeta<C\),
\begin{equation}\label{pfobs:exact-cost}
 \mathcal C_\zeta(g)=v-m+C-\zeta
 =\frac{1-2a-2a^2+(2+a)b}{6(1+a)}-\zeta.
\end{equation}
An attaining chain has \(O(1+\log(1/\zeta))\) edges.
\end{proposition}

\begin{proof}
The values in \eqref{pfobs:profile} match at the joins, and its slopes have
absolute value at most one. On \([1/2,t]\), the difference \(g(x)-ax\)
decreases to zero at \(t\). On \([t,p]\) it is increasing from zero, and on
\([p,1]\) it equals \((1+a)(1-x)\). This proves the lower barrier.
The upper barrier is automatic up to \(t\). On \([t,p]\), its largest
possible violation is at \(p\), and after \(p\) the difference \(g(x)-bx\)
decreases. The upper barrier at \(p\) is equivalent to
$$
 (1-b)p\le t-m.
$$
Substituting the definitions gives exactly
\eqref{pfobs:barrier-condition}. Thus both barriers hold. The global minimum
on \([1/2,1]\) is \(m\), attained at \(t\).

Set
$$
 s=1-C,\qquad \ell=1-2C.
$$
The peak satisfies \(p=1-3C/2\). Moreover,
$$
 \ell-t=\frac{1+2a}{3}(1-t)>0,
 \qquad t<\ell<p<s<1.
$$
Consider any forced-midpoint chain. Let \(n\to q\) be its first edge with
\(n>s\) and \(q\le s\). Every preceding edge lies on the terminal affine
segment of slope \(-1\), so the preceding costs sum to \(1-\zeta-n\).
Admissibility gives
$$
 q\ge2n-1>\ell>t.
$$
If \(q\ge p\), the crossing edge has cost \(n-q\ge n-s\). Therefore the
accumulated cost through this edge is at least \(C-\zeta\).
If \(q<p\), the minimum of the profile on \([q,n]\) is the minimum of its
endpoint values. The crossing cost is consequently
$$
 \bigl(q+n-(1+a+t-m)\bigr)_+
 \ge3(n-s),
$$
where we used \(q\ge2n-1\) and \(3s=2+a+t-m\). Adding the preceding
cost again gives at least \(C-\zeta\).

This crossing edge ends strictly above \(t\). In the remaining part of the
chain, let \(u\to r\) be the first edge with \(u>t\) and \(r\le t\).
Before the required midpoint, its endpoint satisfies \(r\ge1/2\). The
cost of this edge is \(g(r)-m\), since its interval contains the global
tail minimum. Every later edge from \(u'\) to \(r'\) has cost at least
\(g(r')-g(u')\). Telescoping from \(r\) down to \(1/2\) gives another
cost of at least \(v-m\). These are disjoint portions of the chain, so its
total cost is at least \(v-m+C-\zeta\).

For attainment, start at \(1-\zeta\) and double the complementary depth
until reaching \(s\), shortening the final jump if needed. Each edge is
admissible and lies on the affine segment of slope \(-1\), so this part
costs exactly \(C-\zeta\). Continue with
$$
 s\longrightarrow\ell\longrightarrow t\longrightarrow\frac12
 \longrightarrow0.
$$
The first edge is admissible with equality. Its endpoint values coincide,
and the profile is above that value between them, so its cost is zero.
The next edge is admissible because
$$
 2\ell-t\le1
 \quad\Longleftrightarrow\quad
 C\ge\frac{1-t}{4}
 \quad\Longleftrightarrow\quad a\le\frac14.
$$
Its cost is zero because \(t\) is a global tail minimum. The edge from
\(t\) to \(1/2\) is admissible since \(t\le3/4\), and costs \(v-m\).
The final edge costs zero. The initial doubling procedure has logarithmically
many edges. Finally, substitution of \(t=(1+b)/(2(1+a))\) proves the second
identity in \eqref{pfobs:exact-cost}.
\end{proof}

\subsection{The critical packing curve and its verified range}

Let \(a_*\) be the unique root in \((1/5,41/200)\) of
\begin{equation}\label{pfobs:root}
 -48a^3+94a^2+37a-11=0.
\end{equation}
Numerically,
$$
 a_*=0.203243978979\ldots.
$$
For existence, the polynomial in \eqref{pfobs:root} has respective values
\(-28/125\) and \(3807/31250\) at \(1/5\) and \(41/200\). Its derivative
\(-144a^2+188a+37\) is positive on \([0,1/4]\), proving uniqueness in
that interval.

\begin{proposition}[A sharp profile balance on a terminal dimension interval]
\label{pfobs:critical}
For every \(a\in[a_*,1/4]\), set
$$
 b=b_E(a):=\frac{8a^2+8a-1}{a+2}.
$$
There is a profile satisfying both linear barriers and the Lipschitz condition
such that
$$
 \lim_{\zeta\downarrow0}\mathcal C_\zeta(g)=a.
$$
Thus the universal limiting bound \(E(a,b)\) is sharp at
\(b=b_E(a)\) on this parameter interval. In dimension notation this critical
curve is
$$
 D=1+b_E(d-1)=\frac{d(8d-7)}{d+1},
 \qquad 1+a_*\le d\le\frac54.
$$
\end{proposition}

\begin{proof}
We check every hypothesis of Proposition~\ref{pfobs:exact}. Since
\(a\ge a_*>1/5>1/7\),
$$
 b_E(a)-a=\frac{(a+1)(7a-1)}{a+2}>0.
$$
Also
$$
 b_E'(a)=\frac{8a^2+32a+17}{(a+2)^2}>0,
 \qquad b_E(1/4)=\frac23,
$$
so \(a\le b_E(a)\le2/3<1\). The prescribed minimum location simplifies to
$$
 t=\frac{8a+1}{2(a+2)}.
$$
It is increasing in \(a\), and at \(a=1/4\) equals \(2/3\). Hence
\(1/2<t\le2/3<3/4\).
The remaining barrier condition follows from the identity
$$
 (1+b_E(a))^2(1-a)-2(1-b_E(a))(1+a)^2
 =\frac{(1+a)^2}{(a+2)^2}
       \bigl(-48a^3+94a^2+37a-11\bigr).
$$
The right side is nonnegative for \(a\ge a_*\), by the monotonicity just
proved. Proposition~\ref{pfobs:exact} therefore applies. The definition of
\(b_E\) gives
$$
 \frac{1-2a-2a^2+(2+a)b_E(a)}{6(1+a)}=a,
$$
so its exact chain cost is \(a-\zeta\). Taking the limit proves the claim.
The displayed dimension formula follows from
\(8a^2+9a+1=(a+1)(8a+1)\).
\end{proof}

At \(a=6/25\), corresponding to \(d=1.24\), the critical values are
$$
 b=\frac{863}{1400},\quad
 D=\frac{2263}{1400}=1.616428571\ldots,\quad
 t=\frac{73}{112},\quad m=\frac{219}{1400},\quad
 C=\frac{247}{2800}.
$$
At \(a=1/4\), \(b=2/3\), the profile has the particularly simple knots
$$
 (0,0),\quad (1/2,1/3),\quad (2/3,1/6),\quad
 (7/8,3/8),\quad (1,1/4).
$$
Its limiting cost is exactly \(1/4\).

\begin{remark}[Scope of the sharpness statement]
For \(d\in[1+a_*,5/4]\), the critical profile rules out a universal limiting
cost estimate strictly below \(d-1\) at \(D=d(8d-7)/(d+1)\), when only the
present Lipschitz condition, linear barriers, edge costs, and curvature
admissibility are used. It does not rule out a different analytic argument,
additional geometric information, or a proof exploiting a different finite-scale
balance. The examples have not been asserted to be distance-set counterexamples.

The family does not verify sharpness throughout the whole interval where the
packing curve improves on \(2d-1\). At its lower crossing, put
\(a_A=(\sqrt{10}-2)/6\) and \(b_E(a_A)=2a_A\). Then the left side of the
barrier identity above is \(-a_A<0\), so the proposed peak exceeds the upper
barrier. No model optimality claim is made here for
\(a_A\le a<a_*\).
\end{remark}
